\label{sec:introduccion}

La simulación numérica de fenómenos físicos gobernados por ecuaciones diferenciales parciales (PDEs, por sus siglas en inglés) ocupa un lugar central en ingeniería y ciencias aplicadas. En este ámbito, los métodos clásicos, como el método de elementos finitos (FEM, por sus siglas en inglés), el método de volúmenes finitos (FVM, por sus siglas en inglés) y el método de diferencias finitas (FDM, por sus siglas en inglés), ofrecen marcos numéricos consolidados para aproximar soluciones de sistemas diferenciales \cite{grossmann2024}. Aun así, estos métodos pueden volverse costosos ante mallas complejas, formulaciones inversas o escenarios donde se requiere inferencia rápida y el solucionador clásico resulta computacionalmente prohibitivo \cite{karniadakis2021}.

Frente a estas limitaciones, el Aprendizaje Automático Científico (SciML, por sus siglas en inglés) ha surgido como un campo orientado a integrar modelos físicos, datos y métodos de aprendizaje automático \cite{karniadakis2021}. Dentro de este campo, las Redes Neuronales Informadas por Física (PINNs, por sus siglas en inglés), introducidas por Raissi \textit{et al.} \cite{raissi2019}, mostraron que una red neuronal artificial puede aproximar la solución de una PDE e incorporar información física mediante una función de pérdida compuesta por residuos de la ecuación gobernante, condiciones iniciales, condiciones de frontera y datos observados \cite{lagaris1998}. Al usar diferenciación automática y no depender directamente de una malla tradicional, las PINNs también facilitan la formulación de problemas inversos, donde se busca estimar parámetros desconocidos a partir de observaciones parciales \cite{raissi2019}.

Desde 2019, la literatura sobre PINNs ha crecido de forma sostenida. Revisiones recientes documentan aplicaciones en dinámica de fluidos, mecánica estructural, electrofisiología cardíaca y propagación de ondas \cite{karniadakis2021,toscano2025pikans}. Sin embargo, las arquitecturas basadas en perceptrones multicapa (MLP, por sus siglas en inglés) continúan presentando dificultades en problemas con soluciones oscilatorias, fenómenos multiescala, acoplamientos físicos fuertes o componentes de alta frecuencia \cite{karniadakis2021,rahaman2019}.

En este contexto, Liu \textit{et al.} propusieron las Redes de Kolmogorov-Arnold (KAN, por sus siglas en inglés), una arquitectura alternativa inspirada en el Teorema de Representación de Kolmogorov-Arnold, el cual permite expresar funciones multivariadas mediante sumas de funciones univariadas. A partir de este principio, las KAN reemplazan parte de la estructura tradicional de los MLP por funciones aprendibles en las conexiones de la red \cite{liu2024kan}. En un MLP, las funciones de activación suelen ser fijas y se aplican sobre los nodos. En una KAN, en cambio, las funciones univariadas adaptables se aprenden sobre las aristas. Esta diferencia podría favorecer la representación de soluciones con variaciones locales pronunciadas u oscilatorias, sobre todo en escenarios donde las arquitecturas convencionales requieren mayor profundidad o ancho para alcanzar una aproximación adecuada.

Al integrarse con el paradigma de física informada, KAN da lugar a las Redes de Kolmogorov-Arnold Informadas por Física (PIKANs, por sus siglas en inglés). Estas arquitecturas conservan la formulación de pérdida física propia de las PINNs, pero sustituyen la red subyacente por una arquitectura basada en KAN \cite{toscano2025pikans}. Estudios recientes reportan que, en ciertos problemas de referencia, las PIKANs pueden alcanzar desempeños competitivos frente a PINNs convencionales con arquitecturas de menor tamaño \cite{shukla2024fair,wang2025kinn}. No obstante, su comportamiento en problemas oscilatorios, multiescala o asociados a propagación de ondas sigue siendo una oportunidad de investigación, especialmente cuando el interés no se limita al error final, sino a comprender cómo la arquitectura representa distintas regiones del dominio físico \cite{toscano2025pikans}.

De esta manera, las PIKANs constituyen una alternativa arquitectural relevante frente al sesgo espectral y las dificultades de representación observadas en PINNs convencionales. Ahora bien, el panorama actual no debe reducirse a determinar si una arquitectura supera a otra en un conjunto de problemas de referencia. Bajo condiciones adecuadas, ambas pueden alcanzar representaciones de buena calidad. Más allá de esa comparación, la pregunta central de este trabajo se orienta a estudiar cómo incorporar mecanismos que permitan a una arquitectura informada por física identificar, priorizar o focalizar regiones de mayor complejidad durante el entrenamiento. Con ello, se busca mejorar la representación aprendida sin depender únicamente del aumento de profundidad, ancho o número de puntos de colocación \cite{kharazmi2021hpvpinns,lu2021deepxde,shukla2024fair}.